\section{实验总结}

本实验围绕基于颜色与人脸几何结构的人脸区域分割展开，依次完成了实验三的颜色建模、最大连通域填孔与形态学边界检验，以及实验四基于双眼定位的双椭圆空间约束、侧向椭圆候选与低支持放宽阈值重建优化。

在实验三中，原方案采用 $k=1.25$ 的 Cr/Cb 阈值、最大连通域和填孔组合，在 32 张独立测试图上的 F1 为 70.69\%，IoU 为 54.67\%，分别比纯阈值基线提高 8.03 和 9.04 个百分点；附加实验加入 $41\times41$ 膨胀、填孔和腐蚀后，测试集 Recall 提高 4.00 个百分点，但误检增加 19.08\%，总体 F1（70.70\%）基本持平，揭示了单纯依靠颜色与无边界形态学操作无法突破精确率—召回率等值线的瓶颈。

在实验四中，针对“直接从肤色轮廓拟合椭圆易受头发粘连或脸部碎裂带偏”的问题，我们利用灰度黑帽变换与训练集左右眼独立 $3\sigma$ 窗口先验定位双眼（严格与放宽两级面积比回退覆盖 31/32 张验证图，1 张安全回退 Lab3），在以双眼为基准的局部坐标系下构造正中及左右侧向椭圆候选（$\pm 0.16\,d_{\text{eye}}$，偏移门槛 $\Delta S_{\text{lat}}=0.06$）与内外嵌套双椭圆（$\lambda_{\text{in}}=0.82,\;\lambda_{\text{out}}=1.20$）；同时对内椭圆支持率低于 0.18 的样本触发 $k_{\text{rec}}=2.00$ 放宽肤色阈值重选与 $27\times27$ 闭运算受限重建，最后通过内椭圆 $11\times11$ 强闭运算与过渡带 $5\times5$ 弱闭运算完成种子锚定连接。在 32 张验证集上，最终方案将精确率从 64.86\% 提升至 79.88\%（+15.02 个百分点），召回率从 82.70\% 提升至 86.15\%（+3.45 个百分点），全局 F1 从 72.70\% 跃升至 82.89\%（+10.19 个百分点），IoU 从 57.11\% 提升至 70.79\%（+13.68 个百分点），29 张样本提升、仅 2 张微幅波动。
